\uuid{NUXM}
\exo7id{4366}
\titre{exo7 4366}
\auteur{quercia}
\organisation{exo7}
\datecreate{2010-03-12}
\isIndication{false}
\isCorrection{true}
\chapitre{Intégration}
\sousChapitre{Intégrale de Riemann dépendant d'un paramètre}
\cpge{mpsi-mp2i.analyse-asymptotique.a, pcsi.analyse-asymptotique.a, ptsi.analyse-asymptotique.a, mpsi-mp2i.integration.c, pcsi.integration.b, ptsi.integration.b, mp-mpi.calcul-differentiel-optimisation.b, pc.fonctions-plusieurs-variables.c, psi.fonctions-plusieurs-variables.c, bcpst2.integrales-generalisees}
\module{Analyse}
\niveau{L2}
\difficulte{}

\contenu{
\question{
\'Etudier $ \int_{t=0}^{+\infty} \frac{e^{-tx}}{1+t^2}\,d t$.
}
\reponse{
Fonction de $x$ de classe $\mathcal{C}^\infty$ sur $]0,+\infty[$, décroissante
de limite $\pi/2$ en $0^+$ et $0$ en $+\infty$. Demi-tangente verticale en~$0^+$,
\'Equivalente à $1/x$ en $+\infty$ (par IPP). \'Equation différentielle~:
$f(x) + f''(x) = 1/x$.
}
}
